% !TEX root = ../Main.tex
\chapter{平面的法向量}

平面 $\alpha$ 的\textbf{法向量} $\vec n=(x,y,z)$ 是一条与 $\alpha$ 垂直的向量——它跟 $\alpha$ 内的每一条向量都垂直。坐标法解空间几何，第一步往往就是把法向量求出来。

\section{原理：线面垂直的判定}

要让 $\vec n$ 垂直于平面 $\alpha$，只需让它垂直于 $\alpha$ 内\textbf{两条不共线}的向量即可（这正是线面垂直的判定）。设 $\alpha$ 内有
\[
\vec a=(x_1,y_1,z_1),\qquad \vec b=(x_2,y_2,z_2),
\]
则 $\vec n=(x,y,z)$ 满足
\[
\begin{cases}
\vec n\cdot\vec a=x_1x+y_1y+z_1z=0,\\[2pt]
\vec n\cdot\vec b=x_2x+y_2y+z_2z=0.
\end{cases}
\]
两个方程、三个未知数，含一个自由量。\textbf{取一个分量为定值}（常令 $z=1$，或挑一个让其余坐标都是整数的值），代回解出另两个，就得到一个具体的法向量 $\vec n=(\ ,\ ,\ )$。

\section{快速求法}

解方程组每次都要现算，慢。更快的办法是\textbf{叉乘}：平面内两向量 $\vec a=(x_1,y_1,z_1)$、$\vec b=(x_2,y_2,z_2)$，它们的叉乘
\[
\vec n=\vec a\times\vec b=\pr{\,y_1z_2-y_2z_1,\ \ z_1x_2-z_2x_1,\ \ x_1y_2-x_2y_1\,}
\]
天生同时垂直于 $\vec a$ 和 $\vec b$，所以直接就是一个法向量。

记法是\textbf{错位相乘相减}：求某个分量时，把另外两列的坐标交叉相乘再相减（中间那个分量按 $z_1x_2-z_2x_1$ 的顺序写，相当于带一个负号）。

\ex{平面内有 $\vec a=(1,0,1)$、$\vec b=(1,1,0)$，求一个法向量。}

\sol{
    直接叉乘：
    \[
    \vec n=\vec a\times\vec b=\pr{\,0\cdot0-1\cdot1,\ \ 1\cdot1-0\cdot1,\ \ 1\cdot1-1\cdot0\,}=(-1,\,1,\,1).
    \]
    验证 $\vec n\cdot\vec a=-1+0+1=0$、$\vec n\cdot\vec b=-1+1+0=0$，与两向量都垂直，确是法向量。
}
